\appendix
\renewcommand{\chaptermark}[1]{%
	\markboth{\MakeUppercase{APPENDIX\ \thechapter\ --\ #1}}{}
}
\chapter{Radiobiological aspects of ionizing radiations}\label{ch:radiobiology}
The aim of this chapter is to provide an overview of the most important biological effects induced by ionizing radiation and to define the physics parameters employed in the contexts of radiobiology, dosimetry and medical physics.
\section{Biological damages of ionizing radiation}\label{sec:biological_damages_ionizing_radiation}
Both the conventional radiotherapy (which uses X-rays, i.e., photons) and the more advanced Particle Therapy (PT) take advantage of the ionizing radiation to damage the cancerous tissues of a patient, triggering the death of the diseased cells which they consist of. The ionizing radiation produces various effects on biological tissues, which largely depend on the beam composition, its energy and the properties of the irradiated tissues. To beat cancer, it is not necessary to eliminate all the cancerous cells, but it is sufficient to impede their indefinite reproduction through a reproductive cell death. To achieve this goal, ionizing radiation need to kill the cell DNA by means of either direct or indirect damage. Direct damage is the most uncommon one, as the DNA occupies only $2$--$3\%$ of the total cell volume. Direct damage can be divided into Single Strand Breaks (SSB) and Double Strand Breaks (DSB): in SSBs, a break occurs in only one nucleotide strand of the DNA double helix, whereas in DSBs both nucleotide strands are damaged. Both SSB and DSB are illustrated in \autoref{fig:dna_damages}.
\begin{figure}[H]
	\centering
	\includegraphics[width=0.98\linewidth]{danni_dna.jpg}
	\caption[Tracking region]{Schematic representation of SSBs (left), DSBs (center), and clusters of DSBs (right). From \cite{unipv_conference}.}
	\label{fig:dna_damages}
\end{figure}
Indirect damage occurs when the ionizing radiation does not hit the DNA directly, but rather the neighboring cell elements. As water makes $\gtrsim70\%$ of the total cell mass \cite{cooper2000cell}, the ionizing radiation often hits water, leading to the release of free radicals. Free radicals are neutral atoms or molecules possessing one unpaired electron in the outermost orbital, a feature which makes them highly unstable and reactive. The chain of chemical reactions (summarized below) which lead to the formation of free radicals is called radiolysis of water. The byproducts of radiolysis consists of electrons, atoms, molecules, and ions which, by diffusing, indirectly damage DNA, leading to cell death. The ionization of a water molecule \ce{H2O} is given by:
\begin{equation}
	\ce{H2O ->[Ionizing radiation] H2O^{+} + e^-},
	\label{eq:radiolisi1}
\end{equation}
from which a positive ion \ce{H2O^{+}} and an electron \ce{e^-} are produced. After a certain time interval, some traveling electrons lose a significant part of their kinetic energy and are captured by another water molecule, resulting in the formation of a negative ion \ce{H2O^{-}}:
\begin{equation}
	\ce{H2O + e^{-} ->[Electron capture] H2O^{-}}.
	\label{eq:radiolisi2}
\end{equation}
Subsequently, the products of \autoref{eq:radiolisi1} and \autoref{eq:radiolisi2} dissociate as follows:
\begin{subequations}
	\begin{align}
		\label{eq:dissociazione1}
		\ce{H2O^{+} ->[Dissociation] H^{+} + OH^{.}}\\
		\label{eq:dissociazione2}
		\ce{H2O^{-} ->[Dissociation] OH^{-} + H^{.}},
	\end{align}
\end{subequations}
whose products are the hydrogen \ce{H^{+}} and the hydroxide \ce{OH^{-}} ions, along with the hydroxyl \ce{OH^{.}} and the hydrogen \ce{H^{.}} free radicals. Among the various chemical reactions that can subsequently occur (whose reactants are the products of \autoref{eq:dissociazione1} and \autoref{eq:dissociazione2}), the following are reported:
\begin{subequations}
	\begin{align}
		\label{eq:prodotto1}
		\ce{H^{.} + OH^{.} &-> H2O}\\
		\label{eq:prodotto2}
		\ce{H^{+} + OH^{-} &-> H2O}\\
		\label{eq:prodotto3}
		\ce{2OH^{.} &-> H2O2}.
	\end{align}
\end{subequations}
Clearly, \autoref{eq:prodotto1} and \autoref{eq:prodotto2} are harmless reactions (they generate water molecules), whereas \autoref{eq:prodotto3} leads to the formation of hydrogen peroxide \ce{H2O2}, a substance that highly harmful to the cell. Indeed, \ce{H2O2} is capable of altering the structure and function of proteins, mutating nucleotide bases and generating SSBs in DNA \cite{ros}. Generally, \ce{H2O2} and other oxygen-related chemical species (such as the two radicals superoxide anion \ce{O2^{-.}} and hydroxyl \ce{OH^{.}}) are grouped in the wider ROS (Reactive Oxygen Species) family, whose elements can all lead to the reproductive cell death.
\section{Dosimetric quantities}\label{sec:dosimetric_quantities}
In radiobiology, the branch of dosimetry is fundamental to define and quantify the physical quantities that describe the irradiation of matter and the consequent biological effects. Dosimetric studies are particularly important at the treatment planning stage, as they allow to assess the correct balance between the damaging irradiation of the tumor and the risks of toxicity to healthy organs. One of the aims of the International Commission on Radiation Units and Measurements (ICRU) is to define the physics parameters used in dosimetry, some of which are reported in the following.
\subsection{Absorbed, equivalent and effective dose}\label{ssec:dose_definition}
First and foremost, the basis quantity in dosimetry is the particle fluence $\phi$, defined as the number of particles $dN$ (for example, the protons within a beam) passing through an unit area $dA$ perpendicular to the beam direction. The fluence can be expressed in its differential form as follows:
\begin{equation}
	\phi=\frac{dN}{dA}
	\label{eq:fluence}.
\end{equation}
Using the particle fluence, it is possible to define the absorbed dose $D$ as the energy $dE$ absorbed by an irradiated material of unit mass $dm$:
\begin{equation}
	D=\frac{dE}{dm}=\frac{\left(dE/dx\right)\Delta x N}{\rho \Delta x A}=\phi\frac{\left(dE/dx\right)}{\rho},
	\label{eq:dose_as}
\end{equation}
where $dE/dx$ is the stopping power defined in \autoref{eq:bethe_bloch}, $N$ is the number of particles in the beam and $\rho$ is the volumetric mass density of the irradiated material (obtained from its cross-sectional area $A$ and its thickness $\Delta x$, which is parallel to the beam direction). According to the International System of Units (SI), the unit of measurement of the absorbed dose $D$ is the Gray (Gy), and represents the amount of radiative energy received by the patient during therapy divided by the mass of the irradiated tissue.

Generally, when two tissues absorb the same dose, the subsequent biological effects may be different. This is connected to the fact that the biological damage varies depending on the type and energy of the absorbed radiation as well as on repair mechanisms activated by the irradiated tissue after treatment. In radiotherapy, the local biological effects can be described by several quantities, namely the LET (see ?? let), the RBE (see ?? rbe) and the OER (see ?? OER). Due to the inherent complexities of these quantities, everyday radiation protection uses a simpler way to asses the biological effectiveness of a given radiation, which is based on the concept of equivalent dose $H$:
\begin{equation}
	H=\sum_{R}w_RD_{R},
	\label{eq:dose_eq}
\end{equation}
where $D_{R}$ is the absorbed dose relative to the $R$-th radiation (measured in Sievert (Sv) in the SI) and $w_R$ weights the hazard of the $R$-th radiation. $w_R$ is a dimensionless weighting factor given by the ratio between the biological damage produced by the absorption of $1\,\mathrm{Gy}$ of radiation $R$ and the biological damage produced by $1\,\mathrm{Gy}$ of $\gamma$ rays. $H$ is still an incomplete quantity, as it does not quantify biological effects in relation to the type of organ or tissue irradiated. For this reason, the effective dose $E$ is defined (always measured in Sievert), which includes the biological effects due to the type of radiation $R$ and to the specific response of the irradiated tissue $T$:
\begin{equation}
	E=\sum_{T}w_TH=\sum_{T}w_T\sum_{R}w_RD_{R},
	\label{eq:dose_eff}
\end{equation}
where $w_T$ is a weighting factor for the $T$-th tissue such that $\sum_{T}w_T=1$, where the sum runs over all tissues of the human body. $w_T$ is larger for tissues that are more sensitive to radiation: they respond better to the treatment, but they are also more prone to develop radiation-induced cancer.
\subsection{Linear energy transfer}\label{ssec:let}
The Linear Energy Transfer (LET) of a ionizing radiation is the amount of energy $dE$ transferred to the irradiated material per unit length of the track $dx$:
\begin{equation}
	\text{LET}=\frac{dE_\Delta}{dx}.
	\label{eq:let}
\end{equation}
Although the LET has a very similar definition to the stopping power $dE/dx$ defined in \autoref{eq:bethe_bloch}, they are not quite the same quantity. Indeed, the stopping power is related to the energy lost by the particle beam in the medium it crosses (per unit distance), LET refers to the energy absorbed by the medium (per unit distance). In addition, in \autoref{eq:let} a $\Delta$ appears because the definition of LET takes into account only the small energy transfers close to the particle's trajectory, excluding the more energetic ionizations from long-range $\delta$ electrons (whose energy is greater than a threshold $\Delta$). In the mathematical limit $\Delta\rightarrow+\infty$, there would no longer be electrons with $E>\Delta$, meaning that the stopping power and LET would be equivalent quantities. These considerations allow us to derive the relation linking the stopping power and LET:
\begin{equation}
	\text{LET}=\frac{dE_\Delta}{dx}-K_{E>\Delta},
	\label{eq:let&stoppingpower}
\end{equation}
where $K_{E>\Delta}$ is the kinetic energy of electrons with energy $E>\Delta$.

The LET value is inherently connected to particles' ionization density: high-LET radiations are densely ionizing, while low-LET radiations produce sparse ionization events, as illustrated in \autoref{fig:ioniz_density}. As they generate more frequent and clustered ionization events, higher-LET radiations cause greater biological damage to tissues, as schematized in \autoref{fig:damage_let}. \autoref{fig:damage_let} shows that low-LET radiations, such as X-rays or $\gamma$ rays, mainly induce isolated DNA damage leading to SSBs, while high-LET radiations, including $\alpha$ particles or heavy ions, generate more destructive DNA damage, resulting in DSBs \cite{Roobol2020-fp}.
\begin{figure}[H]
	\centering
	\begin{subfigure}[b]{0.48\linewidth}
		\centering
		\includegraphics[width=0.8\linewidth]{ioniz_density.jpg}
		\caption{Variation of ionization density associated with different types of radiation. From \cite{hall2012radiobiology}.}
		\label{fig:ioniz_density}
	\end{subfigure}
	\hfill
	\begin{subfigure}[b]{0.48\linewidth}
		\centering
		\includegraphics[width=\linewidth]{damage_let.png}
		\caption{Typical DNA damage caused by low-LET radiation (top) and high-LET radiation (bottom). From \cite{Fabbrizi_Parsons_2022}.}
		\label{fig:let}
	\end{subfigure}
	\caption[Ionization density and LET]{How different values of LET lead to diverse ionization densities and biological consequences.}
	\label{fig:damage_let}
\end{figure}
From \autoref{fig:proton_carbon_let}, it can be viewed that the LET values of protons and carbons rise with the penetration depth, in accordance with the fact that the stopping power increases when particles slow down. On the other hand, in the entrance channel protons and carbon ions show discrepancies: the former cause damages every $90\,\mathrm{nm}$ (insufficient to hit the DNA double strand, whose diameter is $2\,\mathrm{nm}$), while carbons damage the DNA every $4\,\mathrm{nm}$ (sufficient to generate an SSB). Similarly, at the Bragg peak protons and carbon ions cause damage every $4\,\mathrm{nm}$ and every $0.3\,\mathrm{nm}$ (sufficient to generate SSBs and DSBs, respectively). Therefore, carbon ions are more damaging than protons near the Bragg maximum, but in the entrance channel they cause much more damage than protons. This difference must be taken into account during treatment planning in order to adapt the patient's characteristics to the beam properties.
\begin{figure}[H]
	\centering
	\includegraphics[width=0.9\linewidth]{proton_carbon_let.jpg}
	\caption[LET effects in protons and carbons]{Microscopic distribution of proton and carbon ionizations on a nm scale, along with the corresponding LET, depth reached, kinetic energy and average distance of damage creation \cite{slide_spighi}.}
	\label{fig:proton_carbon_let}
\end{figure}
\subsection{Relative biological effectiveness}\label{ssec:rbe}
The Relative Biological Effectiveness (RBE) of an ionizing radiation $R$ is a dimensionless number given by the ratio between the dose $D_\gamma$ released by $\gamma$ rays and the dose $D_R$ released by a given radiation $R$ needed to achieve the same biological effect, for a fixed cell survival rate $S$:
\begin{equation}
	\text{RBE}=\frac{D^S_\gamma}{D^S_R}.
	\label{eq:rbe}
\end{equation}
For example, if a radiation $R$ has $\text{RBE}>1$, this means that a smaller dose of $R$ is needed to achieve the same biological damage as $\gamma$ rays (in other words, radiation $R$ is more effective than $\gamma$ rays). Evaluating the RBE is rather complex, since it depends on the type of radiation and its energy, the LET (see \autoref{fig:let_rbe}), the cell type and its survival rate (see \autoref{fig:survival_dose}), the oxygenation status (see ?? OER sec) and whether fractionation is employed \cite{Choi2012-ef}. The \autoref{tab:rbe} reports some values of RBE for several radiations.



%for analytical formula of biological dose refer to https://iopscience.iop.org/article/10.1088/1367-2630/10/7/075003/pdf





\chapter{\texorpdfstring{$\chi^2$}{χ²} minimization method}\label{ch:chi_squared_min}
The $\chi^2$ minimization method is founded on the $\chi^2$ test, a probabilistic test that provides criteria for verifying the consistency of theoretical hypotheses with a set of experimental data. Specifically, the $\chi^2$ test is based on the evaluation of the $\chi^2$ function, reported in \autoref{eq:chi_square_method}, which expresses the discrepancy between theory and experimental data. This deviation is calculated using the sum of the squared differences between the measured values $y_k$ and the expected values $f\left(x_k;\alpha_j\right)$, and depends on the measurements $x_k$ and the parameters $\alpha_j$:
\begin{equation}
	\chi^2=\sum_{k=1}^N\left(\frac{y_k-f\left(x_k;\alpha_j\right)}{\sigma_k}\right)^2,
	\label{eq:chi_square_method}
\end{equation}
where $\sigma_k$ is the uncertainty associated with the $N$ measurements $y_k$ \cite{Fornasini2008}. The minimum $\chi^2$ method consists of deriving the parameters\footnote{From a mathematical point of view, the parameters are constraints on the minimization problem.} $\alpha_j$ to minimize the $\chi^2$ function in \autoref{eq:chi_square_method}.

Applying the $\chi^2$ minimization method on the three mass estimates given in \autoref{eq:a1}, \autoref{eq:a2} and \autoref{eq:a3}, the $\chi^2$ function becomes:
\begin{equation}
	\chi^2=\frac{\left(\text{ToF}-T\right)^2}{\sigma^2_{\text{ToF}}}+\frac{\left(p-P\right)^2}{\sigma^2_{p}}+\frac{\left(E_\text{kin}-K\right)^2}{\sigma^2_{E_\text{kin}}}+\transp{\vect{A}}\matr{B}\vect{A},
	\label{eq:chi_square}
\end{equation}
where $\text{ToF}$, $p$ and $E_\text{kin}$ are the physical quantities reconstructed using the FOOT electronic setup, $\sigma_{\text{ToF}}$, $\sigma_{p}$ and $\sigma_{E_\text{kin}}$ the relative uncertainties and $T$, $P$ and $K$ are three of the four minimization parameters. The \autoref{eq:chi_square} contains also the vector $\vect{A}$ (and its transposed $\transp{\vect{A}}$), given by the following expression:
\begin{equation}
	\vect{A}=
	\begin{pmatrix}
		A_1-A\\
		A_2-A \\
		A_3-A\\
	\end{pmatrix}
	,
	\label{eq:vector_A}
\end{equation}
where $A_1$, $A_2$ e $A_3$ are the experimental masses given by \autoref{eq:a1}, \autoref{eq:a2} and \autoref{eq:a3}, and $A$ is the fourth minimization parameter. The matrix $\matr{B}$ contains the uncertainties associated with $A_1$, $A_2$ and $A_3$ estimates, which, as already discussed, are not independent. Indeed, the correlation between the uncertainties of the mass number measurements is encompassed by the correlation matrix $\matr{C}$, which is related to $\matr{B}$ as follows:
\begin{equation}
	\matr{B}=\left(\matr{C}\transp{\matr{C}}\right)^{-1},
	\label{eq:B_funct_C}
\end{equation}
where, specifically, $\matr{C}$ entries are made explicit in \autoref{eq:correlation_matrix} \cite{foot_cdr,valle2019design}:
\begin{equation}
	\matr{C}=
	\begin{pmatrix}
		\frac{\partial A_1}{\partial T}dT&\frac{\partial A_1}{\partial P}dP&0\\
		\frac{\partial A_2}{\partial T}dT&0&\frac{\partial A_2}{\partial K}dK\\
		0&\frac{\partial A_3}{\partial P}dP&\frac{\partial A_3}{\partial K}dK\\
	\end{pmatrix}
	.
	\label{eq:correlation_matrix}
\end{equation}

If the parameters $\left(T,P,K,A\right)$ were thought of as variables of a four-dimensional space, the problem of minimizing $\chi^2$ is analogous to a problem of finding the minimum of a vector-valued function. Therefore, to obtain parameters $\left(T,P,K,A\right)$ that are the closest to the ``true'' parameter values, a minimization fit can be performed. The four-dimensional point where the $\chi^2$ assumes the minimum value will then correspond to the values of $\left(T,P,K,A\right)$ that are the closest to the ``true'' parameter values. A common strategy used to find minima within the ROOT framework is to employ Minuit \cite{JAMES1975343}, which performs multiple minimization steps to find a global $\chi^2$ minimum rather than a local minimum, ensuring that the final output parameters are sufficiently accurate.
\chapter{Augmented Lagrangian Method}\label{ch:alm_squared_min}
The Augmented Lagrangian Method (ALM) is an algorithm that addresses a constrained minimization problem of a multi-variable function $f\left(\vect{x}\right)$ by employing an unconstrained minimization problem in a larger parameter space. The ALM is a ``modified'' Lagrange multiplier method in which the minimization of the Lagrangian function $\mathcal{L}\left(\vect{x};\vect{\lambda}\right)$ with respect to the vector variable $\vect{x}$ is replaced by the minimization of an ``augmented'' Lagrangian $\mathcal{L}_{ALM}\left(\vect{x};\vect{\lambda},\mu\right)$ defined as:
\begin{equation}
	\begin{aligned} 
		\mathcal{L}_{ALM}\left(\vect{x};\vect{\lambda},\mu\right)&=\mathcal{L}\left(\vect{x};\vect{\lambda}\right)+\frac{1}{2\mu}\sum_{i=1}^{m}c_i^2\left(\vect{x}\right)=\\ 
		&=f\left(\vect{x}\right)+\sum_{i=1}^{m}\lambda_ic_i\left(\vect{x}\right)+\frac{1}{2\mu}\sum_{i=1}^{m}c_i^2\left(\vect{x}\right),
	\end{aligned}
	\label{eq:augmented_lagrangian_funct}
\end{equation}
where $f\left(\vect{x}\right)$ is the original function to be minimized, $\vect{\lambda}$ is the $m$-dimensional vector containing the Lagrange multipliers $\lambda_i$ associated with the $m$ constraints $c_i\left(\vect{x}\right)$ and $\mu$ is a positive parameter. The final summation term in \autoref{eq:augmented_lagrangian_funct} is usually referred to as the ``penalty term'', and it is a function of the problem constraints weighted with a penalty parameter $\mu$. If the $c_i\left(\vect{x}\right)$ are zero, the penalty term vanishes, while it is large if $c_i\ne0$ and $\mu\rightarrow0$.

In ALM, successive unconstrained minimizations of \autoref{eq:augmented_lagrangian_funct} are performed: starting from a solution $\left(\vect{x}^0,\vect{\lambda}^0\right)$, given by the unconstrained minimization problem of $f\left(\vect{x}\right)$, the ALM arrives at the $k$-th interaction, obtaining new values $\left(\vect{x}^k,\vect{\lambda}^k\right)$. During these steps, also $\mu$ is modified up to the interaction in which the tolerance level set at the beginning of ALM is reached. At that point, a vector $\overline{\vect{x}}$ is found, representing a point in the parameter space where the $\mathcal{L}\left(\vect{x};\vect{\lambda},\mu\right)$ reaches a minimum value \cite{Cho_2016}.

Applying the ALM method on the three mass estimates given in \autoref{eq:a1}, \autoref{eq:a2} and \autoref{eq:a3}, whose parameters are $\left(T,P,K,A\right)$, it is possible to retrieve the first term of \autoref{eq:augmented_lagrangian_funct}. As shown in \autoref{eq:first_term_ALM}, this term is analogous to that of the \autoref{eq:chi_square}, proving the compatibility of $\chi^2$ and ALM methods:
\begin{equation}
	f\left(T,P,K\right)=\frac{\left(\text{ToF}-T\right)^2}{\sigma^2_{\text{ToF}}}+\frac{\left(p-P\right)^2}{\sigma^2_{p}}+\frac{\left(E_\text{kin}-K\right)^2}{\sigma^2_{E_\text{kin}}}.
	\label{eq:first_term_ALM}
\end{equation}
The last two terms of \autoref{eq:augmented_lagrangian_funct} are given by:
\begin{equation}
	\sum_{i=1}^{3}\lambda_ic_i\left(A\right)+\frac{1}{2\mu}\sum_{i=1}^{3}c_i^2\left(A\right),
	\label{eq:other_terms_ALM}
\end{equation}
where $c_i\left(A\right)=A_i-A$ \cite{valle2019design}. Combining \autoref{eq:first_term_ALM} and \autoref{eq:other_terms_ALM}, it is possible to obtain the following function $\mathcal{L}_{ALM}\left(\vect{x};\vect{\lambda},\mu\right)$:
\begin{equation}
	\begin{aligned} 
		\mathcal{L}_{ALM}\left(T,P,K,A;\vect{\lambda},\mu\right)&=\frac{\left(\text{ToF}-T\right)^2}{\sigma^2_{\text{ToF}}}+\frac{\left(p-P\right)^2}{\sigma^2_{p}}+\frac{\left(E_\text{kin}-K\right)^2}{\sigma^2_{E_\text{kin}}}+\\ 
		&+\lambda_1\left(A_1-A\right)+\lambda_2\left(A_2-A\right)+\lambda_3\left(A_3-A\right)+\\ 
		&+\frac{1}{2\mu}\left(\left(A_1-A\right)^2+\left(A_2-A\right)^2+\left(A_3-A\right)^2\right).
	\end{aligned}
	\label{eq:alm_analysis}
\end{equation}